\documentclass[11pt]{article}
\usepackage[margin=1in]{geometry}
\usepackage{amsmath,amssymb,amsthm,mathtools}
\usepackage{booktabs,longtable,array}
\usepackage[hidelinks]{hyperref}

\newtheorem{definition}{Definition}
\newtheorem{theorem}{Theorem}
\newtheorem{proposition}{Proposition}
\newtheorem{corollary}{Corollary}
\newtheorem{lemma}{Lemma}
\newtheorem{example}{Countermodel}
\newcommand{\R}{\mathbb{R}}
\newcommand{\rank}{\operatorname{rank}}
\newcommand{\im}{\operatorname{im}}
\newcommand{\QS}{Q^{\mathrm{src}}}
\newcommand{\QY}{Q^{\mathrm{out}}}

\title{V22 P2-T01: Local Source Information-Capacity Theorem\\
\large Coordinate- and Gauge-Aware Draft/Control Manuscript}
\author{Ontology Formalizer with same-context Refuter audit}
\date{2026-08-09}

\begin{document}
\maketitle

\begin{abstract}
This draft/control packet proves a conditional local necessary condition for
reconstructing operationally distinguishable candidate geometry from a
finite-order source state.  Source and output redundancies are removed before
rank is evaluated.  If the induced reconstruction factors through a source
statistic of quotient rank at most \(r\), then no completely typed operational
probe family can realize more than \(r\) independent local response directions.
The result recovers the protected P5 factor-through-amplitude rank-one theorem
as a corollary, supplies a controlled stochastic analogue, and classifies
deterministic, stochastic, quotient-observable, partial-domain, marked-domain,
non-smooth, and discrete escape cases.  It constructs no effective metric,
adopts no source structure, and proves no global ontology impossibility.
\end{abstract}

\section{Authority, status, and exact burden}

The active milestone is \texttt{effective\_metric\_g\_eff}.  The burden is:
\emph{prove a coordinate- and gauge-aware source information-capacity necessary
condition and classify escape routes without importing target geometry}.
The source manifold, bundle, jet order, redundancy actions, output space, and
operational probes below are proposal-only mathematical types used to state a
test.  They are not canonical ontology and are not identified with physical
spacetime or an adopted Lorentzian metric.

The P1-T04 physical-typing lock governs every operational label.  Cone
deformation, anisotropic tidal response, polarization response, clock-rate
shift, volume/shear variation, and curvature response are merely candidate
probe slots here.  A slot contributes to operational rank only when its source
observable, source-to-device bridge, scale and units, domain, invariant
readout, and falsification rule are all independently source-derived and
recorded.  Naming a slot supplies none of those witnesses.

\section{Local objects and assumptions}

\begin{definition}[Finite-order source state]
Let \(A\) be a smooth source manifold retained as explicit primitive debt and
let \(E_{\rm src}\to A\) be a proposal-only smooth bundle.  For a finite
\(k\geq 0\), let \(S\) be a finite-dimensional smooth submanifold of a local
finite-order state space \(J^kE_{\rm src}\) on an open source patch.  A point
\(z\in S\) is a local source state.  No target atlas, metric, connection,
measure, or physical clock is part of this definition.
\end{definition}

\begin{definition}[Source and output redundancy]
On a neighborhood \(S_0\) of \(z\), let source relabelings act smoothly with a
locally constant-rank vertical distribution
\[
 V_{\rm src}(z)=\{\text{tangents to the source-redundancy orbit through }z\}
 \subset T_zS .
\]
Let \(Y\) be an abstract candidate-output manifold and let coordinate or
candidate-gauge relabelings have a locally constant-rank vertical distribution
\[
 V_{\rm out}(y)\subset T_yY .
\]
Assume local quotient manifolds \(\bar S\), \(\bar Y\) and submersions
\(\pi_S:S_0\to\bar S\), \(\pi_Y:Y_0\to\bar Y\) exist, with kernels equal to the
vertical distributions.  Define quotient tangent spaces
\[
 \QS_z=T_zS/V_{\rm src}(z)\cong T_{[z]}\bar S,\qquad
 \QY_y=T_yY/V_{\rm out}(y)\cong T_{[y]}\bar Y .
\]
Only this local regular quotient is assumed; a global orbit space need not
exist.
\end{definition}

\begin{definition}[Equivariant reconstruction]
A \(C^1\) map \(R:S_0\to Y_0\) is admissible only when it respects the declared
source and output equivalences.  Equivalently, there is a unique local map
\(\bar R:\bar S\to\bar Y\) satisfying
\[
 \pi_Y\circ R=\bar R\circ\pi_S .
\]
Its derivative \(d\bar R_{[z]}:\QS_z\to\QY_{R(z)}\) is the reconstruction
derivative.  This definition never selects a representative of either orbit.
\end{definition}

\begin{definition}[Source statistic and information capacity]
Suppose the quotient reconstruction factors locally as
\[
 \bar R=F\circ I,\qquad
 I:\bar S\to Z,\quad F:Z\to\bar Y ,
\]
where \(I\) and \(F\) are \(C^1\).  The local source information capacity of
this route at \([z]\) is
\[
 c_I([z])=\rank(dI_{[z]}).
\]
The route has local capacity at most \(r\) when \(c_I\leq r\) on the tested
constant-rank patch.
\end{definition}

\begin{definition}[Admitted operational probe and required response rank]
An admitted probe family is a \(C^1\) invariant map
\(P:\bar Y\to\R^m\) whose components each possess the complete physical-typing
provenance required above.  Let \(W\subseteq\im(d\bar R_{[z]})\) be a candidate
output perturbation subspace.  It is operationally independent for \(P\) when
\(dP_{[y]}|_W\) is injective.  Its required response rank is
\(n=\dim W\).  This number is invariant under changes of source coordinates,
output coordinates, and local quotient charts.  It is not inferred from the
number of symbols used to write a tensor.
\end{definition}

The theorem uses the following explicit assumptions:
\begin{enumerate}
  \item \(S\), \(Y\), and the local quotient manifolds are finite-dimensional
        and \(C^1\) on the tested patch.
  \item Source and output vertical distributions have locally constant rank,
        and the displayed local quotient submersions exist.
  \item \(R\) is \(C^1\) and equivariant, so \(\bar R\) is defined without a
        section or representative choice.
  \item The factorization \(\bar R=F\circ I\) is exact on the tested patch.
  \item Every counted component of \(P\) is invariant and completely
        source-derived under the P1-T04 typing contract.
  \item The claimed output directions lie in \(\im(d\bar R)\); unattained
        directions cannot be credited to the reconstruction.
\end{enumerate}

\section{The quotient-local theorem}

\begin{theorem}[Quotient-local source information-capacity theorem]
\label{thm:capacity}
Under the assumptions above, let \(c_I([z])\leq r\).  For every admitted
operational probe family \(P\),
\[
 \rank\!\left(d(P\circ\bar R)_{[z]}\right)\leq r .
\]
If an adequate reconstruction must realize an \(n\)-dimensional
operationally independent subspace
\(W\subseteq\im(d\bar R_{[z]})\), then necessarily
\[
 n\leq c_I([z])\leq r .
\]
Thus failure of \(n\leq r\) is a scoped local obstruction to that exact
factor-through route, not to all possible source extensions or all global
reconstructions.
\end{theorem}

\begin{proof}
Differentiate the exact quotient factorization:
\[
 d\bar R_{[z]}=dF_{I([z])}\circ dI_{[z]} .
\]
Postcompose with \(dP_{[y]}\).  Rank cannot increase under composition, hence
\[
 \rank d(P\circ\bar R)_{[z]}
 \leq\rank d\bar R_{[z]}
 \leq\rank dI_{[z]}=c_I([z])\leq r .
\]
If \(W\subseteq\im(d\bar R_{[z]})\) and \(dP|_W\) is injective, choose for each
basis vector of \(W\) a preimage in \(\QS_z\).  The composite derivative maps
their span onto an \(n\)-dimensional space, so its rank is at least \(n\).
Combining the inequalities gives \(n\leq c_I([z])\).  All derivatives live on
local quotients.  A change of quotient chart conjugates them by invertible
linear maps and leaves rank unchanged; no orbit representative is selected.
\end{proof}

\begin{corollary}[Protected P5 scalar-amplitude result]
\label{cor:p5}
Let the source state be a patch \(U\), let \(a:U\to\R\), and suppose the
candidate output reconstruction factors as \(R=B\circ a\).  After removing any
declared source and output redundancy, the induced statistic has rank at most
one.  Therefore any completely typed operational requirement with
\(n\geq2\) cannot be realized by this exact factor-through-amplitude route on
the regular patch.
\end{corollary}

\begin{proof}
The codomain of \(a\) is one-dimensional, so \(\rank da\leq1\).  Apply
Theorem~\ref{thm:capacity}.  This reproduces the P5 rank conclusion while
making the quotient, typing, locality, and escape assumptions explicit.
\end{proof}

\section{Stochastic and finite distinguishability forms}

\begin{proposition}[Source-independent stochastic postprocessing]
\label{prop:stochastic}
Let the stochastic output law be a regular dominated \(C^1\) family
\(\mu_{\bar s}=K_{I(\bar s)}\) that depends on the quotient source state only
through \(I:\bar S\to Z\).  Assume the Markov kernel \(K\) has no additional
source-dependent parameter.  Then the local score-tangent span, and hence the
rank of the Fisher information matrix when it exists, is at most
\(\rank dI_{\bar s}\).  Any further source-independent Markov postprocessing
also cannot increase that tangent rank.
\end{proposition}

\begin{proof}
In local coordinates \(\theta=I(\bar s)\), regularity gives the score chain
rule
\[
 d_{\bar s}\log k(\omega\mid I(\bar s))
   =(dI_{\bar s})^{\!*}\,
     d_\theta\log k(\omega\mid\theta).
\]
Every source score therefore lies in the image of \((dI)^{\!*}\), whose
dimension is \(\rank dI\).  The Fisher matrix is a Gram matrix of these scores
and has no larger rank.  Source-independent Markov postprocessing maps tangent
measures linearly, so it cannot enlarge their span.  The proposition makes no
claim about entropy, sample noise, or a kernel carrying a new
source-dependent latent channel.
\end{proof}

\begin{lemma}[Finite distinguishability bound]
\label{lem:finite}
For a deterministic discrete factorization \(S/{\sim_{\rm src}}\to Z\to
Y/{\sim_{\rm out}}\), if the statistic \(Z\) has \(M\) distinguishable values,
then no operational quotient test can realize more than \(M\) distinguishable
output classes.  Hence a requirement for \(N\) classes needs \(M\geq N\).
\end{lemma}

\begin{proof}
Every output class is the image of a statistic value.  A function from an
\(M\)-element image has at most \(M\) distinct images.  This finite statement
is the appropriate bounded check when differentiability is unavailable.
\end{proof}

\section{Operational probe slots}

\begin{longtable}{p{0.20\textwidth}p{0.51\textwidth}p{0.19\textwidth}}
\toprule
Probe slot & Evidence required before its covector counts & Current packet status\\
\midrule
Cone deformation & source observable; propagation/device bridge; invariant
readout; scale; domain; falsifier & conditional slot only\\
Anisotropic tidal response & source response law; extended-body/device bridge;
orientation rule; units; domain; falsifier & conditional slot only\\
Polarization response & source mode definition; detector coupling; quotient
invariant; normalization; domain; falsifier & conditional slot only\\
Clock-rate shift & source clock mechanism; comparison protocol; scale and
units; invariant ratio; domain; falsifier & conditional slot only\\
Volume/shear variation & source volume/shear observable; transport rule;
device bridge; units; domain; falsifier & conditional slot only\\
Curvature response & source-derived curvature surrogate; operational bridge;
gauge-invariant contraction; scale; domain; falsifier & conditional slot only\\
\bottomrule
\end{longtable}

No rank lower bound is instantiated by this list alone.  A later P2-T04
candidate must submit the evidence for each claimed direction, construct the
corresponding invariant differential, and demonstrate independence on the
same quotient patch.

\section{Countermodels and escape classification}

\begin{example}[Raw gauge inflation]
Take source coordinates \((u,\gamma_1,\gamma_2)\) with
\(\gamma_1,\gamma_2\) pure relabeling directions and statistic \(I=u\).
The written source space has dimension three, but
\(\QS\cong\R\) and \(c_I=1\).  Counting the two gauge coordinates would create
a false capacity surplus.
\end{example}

\begin{example}[Representative selection failure]
For the reflection quotient \(x\sim-x\), the invariant coordinate \(q=x^2\)
is regular only away from the fixed orbit.  A formula using a chosen sign of
\(x\) assumes a section that is not supplied by the quotient and cannot be
used in the theorem at the fixed point.
\end{example}

\begin{example}[Two-channel deterministic escape]
For \(I(u,v)=(u,v)\), \(F\) the identity, and two independent invariant probe
covectors, \(c_I=2\).  The rank-two necessity is met.  This is an explicit
added-capacity escape from the scalar corollary, not a derived physical metric.
\end{example}

\begin{example}[Source-independent stochastic noise]
For \(\mu_s=\mathcal N(a(s),1)\), the score is proportional to \(da\) and the
Fisher rank is at most one even if samples are embedded in a high-dimensional
output record.  Noise dimension is not source-parametric capacity.
\end{example}

\begin{example}[Source-dependent latent channel]
For a law with parameter \((u,v)\), the score can have rank two.  The second
source-dependent parameter is explicit new source data and must be charged to
the extension budget; it is not free stochastic postprocessing.
\end{example}

\begin{example}[Partial domain]
The reconstruction \(R(x)=1/x\) on \(x>0\) may satisfy all local hypotheses
there, but gives no result at \(x=0\), on \(x<0\), or about global gluing.
\end{example}

\begin{example}[Marked domain]
Replacing \(a\) by the marked state \((a,m)\) can raise capacity from one to
two when \(dm\) is independent of \(da\).  The mark is an explicit source
extension with provenance and redundancy obligations.
\end{example}

\begin{example}[Rank singularity]
For \(I(u,v)=(u^2,v)\), the rank is two when \(u\neq0\) and one at \(u=0\).
A theorem statement spanning both regions without a stratified or
constant-rank treatment overreaches.
\end{example}

\begin{example}[Discrete injection outside the differential theorem]
A four-class discrete source statistic can inject into four operational
classes even though no derivative exists.  Lemma~\ref{lem:finite}, not
Theorem~\ref{thm:capacity}, tests this case.
\end{example}

\begin{example}[Quotient-observable restriction]
Let \(Y=\R^2\) with translations of the second coordinate declared gauge.
Only the first coordinate descends to \(\bar Y\).  Two raw coordinate
responses therefore represent one quotient-observable response direction.
\end{example}

\begin{center}
\begin{tabular}{p{0.25\textwidth}p{0.29\textwidth}p{0.32\textwidth}}
\toprule
Case & Theorem disposition & Honest escape or stop rule\\
\midrule
Deterministic smooth & Theorem~\ref{thm:capacity} applies & add independent
source channels or abandon exact factorization\\
Stochastic regular & Proposition~\ref{prop:stochastic} applies & declare a
source-dependent channel or nonregular model\\
Quotient observable & count only descended probes & derive additional
invariants, never count gauge coordinates\\
Partial domain & conclusion is patch-local & prove extension and gluing
separately\\
Marked domain & marks enlarge the source state & charge marks as explicit
source data with provenance\\
Non-smooth/discrete & differential theorem is not applicable & use stratified,
topological, or finite distinguishability tests\\
\bottomrule
\end{tabular}
\end{center}

\section{P2-T04 candidate test interface}

A later candidate is testable only if it supplies:
\begin{enumerate}
  \item its source state \(S\), finite order or an explicit nonlocal
        replacement, domain, and regularity;
  \item source redundancy and the quotient tangent or a precise alternative;
  \item abstract output state, output redundancy, and invariant operational
        probe definitions;
  \item the reconstruction and any factor statistic \(I\);
  \item a computed local capacity rank or a declared non-differential test;
  \item complete physical typing and a proof of independence for every counted
        response direction;
  \item treatment of singular strata, partial domains, marks, latent
        variables, randomness, and global gluing; and
  \item a verdict from the vocabulary
        \emph{necessary condition met}, \emph{scoped obstruction},
        \emph{outside differential scope}, or \emph{insufficient evidence}.
\end{enumerate}
Meeting the necessary condition is not sufficient for geometry reconstruction,
physical interpretation, robustness, uniqueness, or Gate B.

\section{Expanded Distance-to-GR matrix}

\begin{longtable}{p{0.27\textwidth}p{0.62\textwidth}}
\toprule
Burden & Status after P2-T01\\
\midrule
Source ontology primitives & unchanged; the smooth source manifold, bundle,
jet order, and local quotient types remain primitive debt or proposal-only\\
Source equivalence \(\mathrm{EqSrc}\) & unchanged; a local redundancy relation
is a theorem assumption, not an adopted canonical equivalence\\
\(\mathrm{RetainH}\) & unchanged; no retention law is derived\\
\(\mathrm{GenH}\) & unchanged; no generator law is derived\\
\(\mathrm{ObsLoc}_{\rm lc}\) & narrowed only by a fail-closed typing interface;
no operational witness is supplied\\
\(\mathrm{Resp}_{\rm lc}\) & unchanged; no universal physical response law is
derived\\
\(M_{\rm src}\) & blocked; the source patch is not physical spacetime\\
\(g_{\rm eff}\) & active burden; a necessary adequacy test is proved, but no
candidate metric is constructed or adopted\\
Matter coupling & unchanged; no universal target coupling is derived\\
Einstein equations & blocked; no target action or field equation is supplied\\
Finite-variation robustness & unchanged; no candidate is stress-tested under
finite variation\\
Benchmark promotion & blocked; no benchmark is executed or passed\\
Gate Chair status & unchanged; no protected Gate review is requested or issued\\
Current route freeze or hard-fail & not triggered; P2-T03 is a distinct
pre-registration package and P2-T04 remains dependent on it\\
\bottomrule
\end{longtable}

The physical Distance-to-GR delta is zero.  This packet discharges the
plan-level obligation to state the necessary-condition theorem but does not
discharge the metric-law burden itself.

\section{Internal review and forbidden overreads}

The Physicist-Mathematician and Refuter perspectives are same-context internal
AI review only.  They are not external mathematical review, independent
replication, proof authority, or permission to use the theorem as a hard
candidate filter.  P2-T02 remains separately human/external-action gated.

The following conclusions are forbidden:
\begin{itemize}
  \item that a raw coordinate dimension or a familiar physical word proves an
        operational lower bound;
  \item that a quotient representative or global quotient exists without the
        stated hypotheses;
  \item that stochastic noise creates source-parametric information;
  \item that the scalar-amplitude corollary rules out added fields, marks,
        non-smooth encodings, discrete refinements, or all future source laws;
  \item that the local theorem constructs \(g_{\rm eff}\), universal coupling,
        Einstein equations, a benchmark pass, or a completed derivation; or
  \item that validation, compilation, checkpointing, or internal agreement
        promotes a scientific claim.
\end{itemize}

After checkpoint, P2-T03 may separately pre-register the source-extension
budget and hard-fail protocol.  P2-T03 is not executed here.

\end{document}
